\chapter{2007:Gerhard Ertl}

\label{chap:chemistry2007}

The Nobel Prize in Chemistry for the year 2007 was awarded to Gerhard Ertl.

\textbf{Motivation:}

for his studies of chemical processes on solid surfaces

\section{Gerhard Ertl}

\subsection*{Early Life and Education}

Gerhard Ertl was born in 1936 in Stuttgart, Germany.

\subsection*{Career and Major Research}

Ertl studied physics at the University of Stuttgart and received his doctorate from the Technical University of Munich in 1965. He became director of the Fritz Haber Institute of the Max Planck Society in Berlin. He established the field of surface chemistry and worked out the elementary steps of reactions on metal surfaces, notably the mechanism of ammonia synthesis on iron and the oxidation of carbon monoxide on platinum.

\subsection*{Nobel Prize-winning Work}

The work for which Gerhard Ertl was awarded the Nobel Prize in Chemistry in 2007 was described by the Nobel Committee as: "for his studies of chemical processes on solid surfaces".

Ertl showed at the level of individual atoms how gases react on well-defined crystal surfaces, identifying the separate steps of catalytic reactions such as adsorption, reaction, and desorption.

\subsection*{Significance and Impact}

His work turned heterogeneous catalysis from an empirical art into a molecular science. The understanding of the Haber--Bosch ammonia synthesis he provided has guided the design of the industrial catalysts used to make fertilizer.

\subsection*{Later Career and Honors}

Ertl headed the physical chemistry department of the Fritz Haber Institute and later served as its director, becoming director emeritus. He received numerous honors, including the Wolf Prize in Chemistry in 1998.

\subsection*{Legacy}

Ertl's surface science methods are now standard tools in catalysis research, and his studies are the reference point for understanding how catalysts work at the atomic level.

\subsection*{Modern Developments}

Ertl's surface chemistry studies, including his visualization of the Haber–Bosch ammonia synthesis at the atomic level, established modern surface science and heterogeneous catalysis. His methods are now used to optimize industrial catalysts for fertilizers, fuel cells, and pollution control, and his work underpins the design of catalytic converters and clean-energy technologies.

\subsection*{Selected Publications}

"Reactions at Surfaces: From Atoms to Complexity" (Angewandte Chemie International Edition, 2008)

\clearpage

\section*{Chapter Summary}

The 2007 prize recognized the founding of modern surface chemistry by Gerhard Ertl, who described the elementary steps of reactions on solid surfaces. This work provided the molecular basis for understanding heterogeneous catalysis, including the industrial synthesis of ammonia.
